\documentclass[11pt]{article}
\usepackage[margin=1in]{geometry}
\usepackage[T1]{fontenc}
\usepackage{lmodern,microtype,hyperref,enumitem,amsmath,amssymb,xcolor,booktabs}
\usepackage[numbers,sort&compress]{natbib}
\hypersetup{colorlinks=true,urlcolor=blue,citecolor=blue,linkcolor=blue}
\title{Q18: Discussion Questions}
\author{DLS 2026 --- Kukla}
\date{October 2026}
\begin{document}
\maketitle
\begin{enumerate}[leftmargin=1.5em]
\item Model activations predict responses in sensory cortex but much less so in high-level reasoning areas. What does this asymmetry suggest about where the brain analogy is load-bearing?
\item Backpropagation is biologically implausible, yet it works spectacularly. Does that make the ``brains learn differently'' objection scientifically interesting or practically irrelevant? Argue a side.
\item If an LLM matches human behavior on a reasoning benchmark but by a different mechanism, in what concrete situations would the mechanism difference actually matter?
\item Neuromorphic hardware promises brain-like energy efficiency. What would have to be true for it to beat GPUs on a real task, not just on power?
\item Which single feature of biological computation, if we could import it faithfully, would you bet most improves current AI---and why?
\end{enumerate}
\end{document}
